% Adaptación del preámbulo de Paper III v1.5 español; no plantilla standalone de Pandoc.
\documentclass[spanish,11pt,a4paper]{article}
\usepackage[margin=1in]{geometry}
\usepackage{fontspec}
\setmainfont{lmroman10-regular.otf}[BoldFont=lmroman10-bold.otf,ItalicFont=lmroman10-italic.otf,BoldItalicFont=lmroman10-bolditalic.otf]
\setsansfont{lmsans10-regular.otf}[BoldFont=lmsans10-bold.otf]
\setmonofont{lmmono10-regular.otf}[Scale=MatchLowercase]
\newlength{\LeanCharWidth}
\usepackage{amsmath,amssymb}
\usepackage[spanish,es-noshorthands,es-tabla]{babel}
\usepackage{microtype,parskip}
\usepackage{graphicx,float,caption}
\usepackage{longtable,booktabs,array,calc}
\usepackage{xcolor}
\usepackage{xurl}
\usepackage{hyperref,bookmark}
\usepackage{etoolbox,needspace}
\setcounter{secnumdepth}{-\maxdimen}
\setlength{\emergencystretch}{3em}
\providecommand{\tightlist}{\setlength{\itemsep}{0pt}\setlength{\parskip}{0pt}}
\newcounter{none}
\AtBeginEnvironment{longtable}{\Needspace{8\baselineskip}\small}
\pretocmd{\section}{\Needspace{5\baselineskip}}{}{}
\pretocmd{\subsection}{\Needspace{4\baselineskip}}{}{}
\pretocmd{\subsubsection}{\Needspace{4\baselineskip}}{}{}
\captionsetup{font=small,labelfont=bf}
\urlstyle{same}
\hypersetup{pdfauthor={Juan Pablo Traverso Gianini},pdflang={es-ES},colorlinks=true,linkcolor=blue,urlcolor=blue,citecolor=blue}
\title{Particiones óptimas en cliques\\bajo perturbaciones locales}
\author{Juan Pablo Traverso Gianini\\\small Investigador independiente, Santiago, Chile\\\small \href{mailto:jtraverso@gmail.com}{jtraverso@gmail.com}\\\small ORCID: 0009-0003-6068-4096}
\date{}\begin{document}\raggedbottom\maketitle
\textbf{Artículo VI de la serie}\\
\textbf{Preprint 1.0}\\
\textbf{Fecha:} 18 de octubre de 2026.

Copyright © 2026 Juan Pablo Traverso Gianini. Licencia CC BY-NC 4.0.

\subsection{Resumen}\label{resumen}

Determinamos exactamente el número de partición en cliques en un entorno
explícito de un grafo completo-split. Escribamos
\(h=\Delta(G\triangle(K_P\vee I_O))\), donde
\(|P|=p=\lfloor(n+1)/3\rfloor\), y sean \(u,m,s\) los números de aristas
del núcleo, exteriores y cruzadas. Si \(p\ge18h+1\), entonces
\(\operatorname{cp}(G)=s-u-2m\), con un óptimo que utiliza bloques de
orden a lo sumo cuatro; al restringirse a triángulos, el óptimo aumenta
exactamente en \(m\). Si \(p\ge63h\), se puede prescribir de manera
independiente que cada arista exterior pertenezca a un bloque de orden
cuatro o cinco. Para \(p\ge16\), cada prescripción admite al menos
\(p^{\lfloor p^2/64\rfloor}\) particiones óptimas. En la región local,
el núcleo de referencia se puede recuperar a partir de los grados.
También determinamos el coste adicional exacto causado por un vértice
excepcional. Este coste puede ser lineal incluso después de optimizar la
cota inferior asociada a la bipartición sobre todos los núcleos. Las
construcciones alcanzan cotas duales explícitas con signo mediante
elecciones de carga acotada y coloreo de aristas con listas; el teorema
de Vizing proporciona el paso de conteo. La cordalidad da una fórmula
exacta de estabilidad local, pero no se supone en los resultados
principales de perturbación.

Para joins completos con \(q\ge2p-1\), dos levantamientos legales
formados solo por \(K_5\) que difieren en \(k\) asignaciones admiten
completaciones óptimas cuya diferencia simétrica de bloques es a lo sumo
\(68k\). Si \(q\in\{2p-1,2p,2p+1\}\) y \(p\ge63\Delta(H)\), los óptimos
con levantamientos formados solo por \(K_5\) y con carga del núcleo a lo
sumo \(7\Delta(H)\) se conectan mediante a lo sumo \(\binom p2-2m\)
pasos de distancia entre bloques a lo sumo \(884\).

\textbf{Palabras clave:} partición en cliques; grafo completo-split;
perturbación local; coloreo de aristas con listas; grafo cordal;
estabilidad exacta.

\subsection{1. Introducción}\label{introducciuxf3n}

Una partición en cliques de un grafo es una partición de su conjunto de
aristas en subgrafos completos. Los bloques pueden compartir vértices y
no tienen que ser cliques maximales. Escribimos \(\operatorname{cp}(G)\)
para el número mínimo de bloques y \(\operatorname{cp}_{\le r}(G)\)
cuando sus órdenes se restringen a lo sumo a \(r\).

Paper IV {[}IV{]} demuestra la cota para grafos cordales de Erdős 81,
con la cota entera exacta para órdenes suficientemente grandes. Este
artículo continúa el estudio de las particiones en cliques más allá de
esa cota global: determina valores exactos bajo perturbaciones locales,
permite prescribir bloques exteriores y estudia la estabilidad y la
conectividad restringida de particiones óptimas. Estos resultados no son
insumos necesarios para completar la demostración de Paper IV.

Los grafos completo-split constituyen una familia de referencia útil
para este problema. Una arista de su parte completa puede servir de base
a un triángulo con un vértice exterior. Si estos triángulos se eligen de
forma compatible, reducen el número de bloques de una sola arista.
Añadir una arista exterior abre otras posibilidades: se puede unir a una
arista del núcleo para formar un \(K_4\), o a un triángulo del núcleo
para formar un \(K_5\). La dificultad consiste en realizar todas estas
elecciones sin agotar las aristas cruzadas necesarias para completar el
núcleo residual.

Demostramos que una cota sobre el número de modificaciones en cada
vértice basta para efectuar esta construcción y alcanzar exactamente una
cota inferior clásica. No es necesario que el grafo sea cordal. La
cordalidad se utiliza solo después, cuando obliga a que el núcleo siga
siendo completo.

\subsubsection{1.1. El modelo local}\label{el-modelo-local}

Sea \(n\ge2\), y pongamos \[
p=\left\lfloor\frac{n+1}{3}\right\rfloor,\qquad
q=n-p,\qquad M(n)=\left\lfloor\frac{n(n+1)}6\right\rfloor.
\] Fijemos una bipartición etiquetada \(V=P\sqcup O\), con \(|P|=p\) y
\(|O|=q\), y sea \(S=K_P\vee I_O\). Entonces \[
q\in\{2p-1,2p,2p+1\},\qquad q\ge p,\qquad
pq-\binom p2=M(n). \tag{1}
\] Para un grafo simple \(G\) sobre este conjunto de vértices,
escribamos \[
\begin{aligned}
C&=\binom p2-e(G[P]), &u&=e(G[P])=\binom p2-C,\\
A&=pq-e_G(P,O), &s&=e_G(P,O)=pq-A,\\
H&=G[O], &m&=e(H).
\end{aligned}
\] Así, \(C\) y \(A\) cuentan las aristas ausentes del núcleo y las
cruzadas, mientras que \(m\) cuenta las aristas exteriores. El parámetro
de modificación local es \[
h=\Delta(G\triangle S). \tag{2}
\] Todas las distancias y todos los tipos que siguen se refieren a esta
bipartición fija. Un clique tiene tipo \((a,b)\) si contiene \(a\)
vértices del núcleo y \(b\) vértices exteriores.

Para \(a\in P\), sean \(c_a\) y \(\alpha_a\) los números de sus vecinos
ausentes dentro del núcleo y en el exterior. Para \(x\in O\), sea
\(\beta_x\) el número de sus vecinos ausentes del núcleo. La definición
de \(h\) da \[
c_a+\alpha_a\le h,\qquad \beta_x+d_H(x)\le h,
\qquad m\le qh/2. \tag{3}
\] Las cotas conjuntas de (3) se utilizan en los pasos de selección.

\Needspace{20\baselineskip}
\subsubsection{1.2. Resultados
principales}\label{resultados-principales}

\Needspace{14\baselineskip}
\textbf{Teorema A (fórmula local exacta).} Si \[
p\ge18h+1, \tag{4}
\] entonces \[
\operatorname{cp}(G)=\operatorname{cp}_{\le4}(G)
=s-u-2m=M(n)+C-A-2m. \tag{5}
\] Existe una partición óptima con \(m\) bloques \(K_4\) de tipo
\((2,2)\), \(u-m\) triángulos de tipo \((2,1)\) y \(s-2u-2m\) bloques de
una sola arista cruzada. Cuando \(q=2p\), basta la condición más débil
\(p\ge18h\). La hipótesis también implica \(m\le u\): por (3),
\(2(u-m)\ge p(p-1)-(p+q)h\ge p(p-1)-(3p+1)h\ge0\). Por tanto, la cota
inferior que aquí se alcanza es exactamente el caso correspondiente de
la cota de Erdős--Faudree--Ordman.

La fórmula contabiliza por separado los tres tipos de modificación. En
esta región, eliminar una arista del núcleo aumenta en uno la expresión
de referencia, eliminar una arista cruzada la reduce en uno y añadir una
arista exterior la reduce en dos. Las constantes de las hipótesis
locales son constantes suficientes obtenidas de la construcción; no se
afirma que sean óptimas.

\textbf{Teorema B (restricción a triángulos).} Bajo (4), \[
\operatorname{cp}_{\le3}(G)=s-u-m=M(n)+C-A-m, \tag{6}
\] y, por tanto, \[
\operatorname{cp}_{\le3}(G)-\operatorname{cp}(G)=m. \tag{7}
\]

Así, incluso una construcción triangular óptima requiere un bloque
adicional por cada arista exterior respecto del óptimo sin restricción.
Se comparan aquí dos familias de bloques permitidos. En el problema de
empaquetamiento de triángulos, tanto el óptimo entero como el fraccional
son \(u+m\), como se demuestra en la Sección 4.

\textbf{Teorema C (bloques exteriores prescritos).} Si \[
p\ge63h, \tag{8}
\] entonces, para toda función \(\sigma:E(H)\to\{4,5\}\), existe una
partición óptima en la que cada arista exterior \(xy\) pertenece a un
bloque \(K_{\sigma(xy)}\) que contiene exactamente dos vértices
exteriores. Todos los demás bloques son triángulos de tipo \((2,1)\) o
aristas cruzadas individuales. Si \(k=|\sigma^{-1}(5)|\), los números de
bloques son \[
N_5=k,\quad N_4=m-k,\quad N_3=u-m-2k,
\quad N_2=s-2u-2m+2k. \tag{9}
\] En particular, \(G\) tiene al menos \(2^m\) particiones óptimas en
cliques distintas.

La prescripción puede variar de una arista a otra. El teorema construye
una partición para cada prescripción. La Sección 9 distingue esta
libertad del cambio acotado disponible en una capa de levantamiento
determinada.

Para joins completos, el caso formado solo por \(K_5\) permite una
conclusión adicional sobre las propias particiones óptimas. La Sección
10 extiende cualquier prefijo terminal legal, controla la partición
completa después de cambiar el levantamiento y conecta mediante pasos
acotados los óptimos de carga acotada con levantamientos formados solo
por \(K_5\). Estos enunciados conservan las hipótesis de join completo y
de usar solo \(K_5\); no se refieren a toda la clase de prescripciones
mixtas.

\subsubsection{1.3. Estrategia de demostración y trabajos
anteriores}\label{estrategia-de-demostraciuxf3n-y-trabajos-anteriores}

Erdős, Faudree y Ordman {[}EFO88{]} demostraron que una bipartición con
\(u,v\) aristas internas y \(s\) aristas cruzadas satisface
\(\operatorname{cp}(G)\ge s-u-v-\min\{u,v\}\). También identificaron la
eficiencia de los bloques \(K_4\) y \(K_5\) con dos vértices en un mismo
lado. Utilizamos una forma asimétrica de su cota cuyos tipos de holgura
cero son precisamente \((1,1),(2,1),(2,2),(3,2)\). La contribución que
aquí se considera es una construcción uniforme que alcanza esa cota en
todo el entorno local indicado, incluso con aristas ausentes del núcleo
y entre las partes, y con prescripciones arbitrarias para las aristas
exteriores.

La construcción tiene dos etapas. Primero, cada arista exterior se
extiende a un clique adecuado. Una cota sobre el número de extensiones
que utilizan cada vértice del núcleo mantiene disponibles suficientes
aristas cruzadas. Después, el núcleo residual se descompone en
triángulos eligiendo un vértice exterior para cada arista residual. Este
es un problema de coloreo de aristas con listas: la lista de una arista
del núcleo está formada por los vértices exteriores cuyas dos aristas
cruzadas siguen disponibles. El teorema de Häggkvist y Janssen
{[}HJ97{]} completa este paso de manera exacta.

El caso de join con un factor independiente tiene un antecedente clásico
más fuerte. Pullman y Donald {[}PD81, Teorema 3{]}, citados aquí
mediante {[}EFO88, p.~94{]} y {[}WZ93, Corolario 3.3{]}, demostraron que
\[
\operatorname{cp}(F\vee I_q)=q|V(F)|-e(F)
\quad\text{si }q\ge\chi'(F),
\] y describieron sus particiones óptimas mediante aristas y triángulos.
El factor \(I_q\) es independiente. Esto incluye nuestro caso \(A=m=0\),
incluso con un grafo arbitrario \(F=G[P]\) en el núcleo. No cubre la
presencia simultánea de aristas cruzadas ausentes y aristas exteriores
añadidas. Wallis y Zhang {[}WZ93, Proposición 3.2{]} dan una
caracterización exacta del número de partición en cliques de un join con
factor independiente mediante una optimización de coloreo de cliques
generalizado. Su Corolario 3.3 da el criterio necesario y suficiente
para la igualdad anterior y lo atribuye a Pullman y Donald. También
recupera un caso de frontera del Teorema D: cuando el vértice
excepcional es adyacente a todo vértice exterior ordinario, trasladarlo
al núcleo produce un join con factor independiente al que se aplica el
criterio clásico. Son antecedentes exactos en las subfamilias de joins;
no proporcionan la construcción con modificaciones simultáneas ni los
levantamientos exteriores prescritos considerados aquí.

Erdős, Ordman y Zalcstein {[}EOZ93{]} plantearon el problema de
partición en cliques para grafos cordales y exhibieron ejemplos
extremales completo-split. Chen, Erdős y Ordman {[}CEO94, Teoremas
1--2{]} demostraron las cotas \(3n^2/16+O(n)\) para grafos split y
\(n^2/6+O(n)\) para grafos completo-split. Su uso de matchings del
núcleo, completaciones triangulares y cuentas de aristas cruzadas
ausentes es un antecedente del método de construcción. Nuestro enunciado
es una igualdad finita bajo las hipótesis locales indicadas, y no una
cota para todo grafo split. En la Sección 8, la familia excepcional es
split cuando \(a=p\); la familia general también incluye grafos no
cordales. Wallis y Wu {[}WW91{]} demostraron la NP-completitud del
problema general para grafos split, lo que es compatible con el
resultado de reconocimiento restringido de este trabajo. La subfamilia
\(a=p\) también es de umbral. Las cotas de Ma por clases de grado
{[}Ma88, Teorema 2{]} ya alcanzan el valor entero en una subregión de
esa familia: cuando el núcleo no está vacío, el vértice excepcional
tiene más vecinos exteriores que el tamaño del núcleo y al menos ese
tamaño de no vecinos exteriores. El Teorema D trata la familia
excepcional completa y determina también sus óptimos fraccional y
triangular. Esta comparación usa el artículo completo en chino, y no
sólo su resumen.

En Okechukwu {[}Oke26, Lemas 4.1--4.3{]} aparecen construcciones
terminales relacionadas. Su Lema 4.2 da una cota superior de partición
triangular bajo sus propias hipótesis de coloreo, capacidad y defecto.
En los subcasos con enlaces completos que cumplen esas hipótesis, la
cota superior coincide con la cota inferior elemental usada en B. Esta
comparación no cubre toda la ventana de modificaciones simultáneas de B.
Ning {[}Nin26, Proposición 3.1{]} alcanza la cota EFO mediante \(K_4\)
para joins de dos uniones especificadas de cliques iguales, bajo sus
condiciones de paridad y tamaño. Monson {[}Mon96{]} estudió los efectos
generales de las eliminaciones sobre el número de partición en cliques.
Las cotas inferiores, los tipos eficientes de bloques, la fórmula de
join con factor independiente y el método de matchings del núcleo son
herramientas clásicas.

El conteo de particiones mínimas de aristas en cliques también tiene
antecedentes clásicos: McGuinness y Rees {[}MR90{]} las enumeran para
grafos línea. Más generalmente, Kwan, Sah y Sawhney {[}KSS23{]} cuentan
descomposiciones de grafos completos en cliques sin la restricción de
número mínimo de bloques usada aquí. El Teorema E da, en cambio, una
cota inferior uniforme para cada patrón prescrito de bloques exteriores
de nuestra región local. La reconfiguración del coloreo de aristas con
listas tiene su propia teoría establecida {[}IKD12{]}. La Sección 10
controla particiones completas en cliques dentro de una clase de carga
acotada con levantamientos sólo de K5 y cambios acotados de bloques; no
es un resultado general de recoloreo de una arista por vez. Estas
comparaciones delimitan las aportaciones candidatas y no constituyen una
afirmación exhaustiva de prioridad histórica.

La Sección 2 presenta el certificado de cota inferior y la herramienta
de completación. Las Secciones 3--5 demuestran A, B y C. Las Secciones
6--7 deducen estabilidad cordal local y optimalidad fraccional. La
Sección 8 recupera el núcleo y determina el defecto del vértice
excepcional; la Sección 9 enuncia el resultado de conteo. La Sección 10
demuestra completación estable y conectividad restringida para joins
completos con levantamientos formados solo por \(K_5\). Los apéndices
presentan una construcción elemental para joins, una localización de los
bloques ineficientes y la demostración del conteo.

\subsection{2. Cotas inferiores y completación
exacta}\label{cotas-inferiores-y-completaciuxf3n-exacta}

\subsubsection{2.1. Una identidad de
holgura}\label{una-identidad-de-holgura}

Para un clique de tipo \((a,b)\), con \(a+b\ge2\), definamos \[
\ell(a,b)=1-ab+\binom a2+2\binom b2
=\frac{(a-b)(a-b-1)+(b-1)(b-2)}2\ge0. \tag{10}
\] La desigualdad se sigue de que el producto de dos enteros
consecutivos es no negativo. Para cualquier partición en cliques
\(\Pi\), contar una vez cada arista presente da \[
|\Pi|-(s-u-2m)=\sum_{Q\in\Pi}\ell(a_Q,b_Q). \tag{11}
\] En consecuencia, \(\operatorname{cp}(G)\ge s-u-2m\) para todo grafo
con la bipartición fija, sin restricción sobre \(h\). Esto también se
deduce de {[}EFO88{]} y coincide con su cota cuando \(m\le u\).

La identidad cuenta de dos maneras las mismas aristas físicas. Cada
arista cruzada aparece una vez en el término que contiene \(ab\), cada
arista del núcleo una vez en el término que contiene \(\binom a2\) y
cada arista exterior una vez en el término que contiene \(\binom b2\).
La contribución restante es una unidad por bloque. Esta cota inferior no
utiliza la elección de ninguna construcción particular.

Ambos productos de (10) se anulan exactamente para \[
(a,b)\in\{(1,1),(2,1),(2,2),(3,2)\}. \tag{12}
\] Por tanto, basta construir una partición que utilice únicamente estos
tipos. Una vez construida tal partición, la identidad también
identificará todos los tipos óptimos.

\subsubsection{2.2. Pares y triángulos disponibles en el
núcleo}\label{pares-y-triuxe1ngulos-disponibles-en-el-nuxfacleo}

Sea \(B\) un grafo de pares prohibidos del núcleo, con
\(\Delta(B)\le D\). Todo conjunto de candidatos con al menos \(D+2\)
vértices contiene un par disponible. En efecto, un vértice prohíbe a lo
sumo \(D\) de los demás. Un conjunto de candidatos con al menos \(2D+3\)
vértices contiene un triángulo disponible: se elige un par disponible y
luego un tercer vértice fuera de sus dos extremos y de sus vecindarios
prohibidos. Se han excluido a lo sumo \(2D+2\) vértices.

\subsubsection{2.3. Completación mediante coloreo de aristas con
listas}\label{completaciuxf3n-mediante-coloreo-de-aristas-con-listas}

Häggkvist y Janssen {[}HJ97{]} demostraron \[
\chi'_{\ell}(K_p)\le p. \tag{13}
\] Aplicamos este resultado a un grafo residual \(J\subseteq G[P]\) del
núcleo. Para cada \(a\in P\), sea \(F_a\subseteq O\) el conjunto de
vértices exteriores cuya arista cruzada hacia \(a\) está ausente o ya se
ha utilizado. A cada \(ab\in E(J)\) se le asigna la lista
\(L_{ab}=O\setminus(F_a\cup F_b)\).

Si todas las listas tienen al menos \(p\) elementos, se extienden las
listas al grafo completo abstracto sobre \(P\) asignando \(O\) a cada
par no residual. Estas listas añadidas tienen tamaño \(q\ge p\), de modo
que se aplica (13). Se restringe el coloreo propio con listas resultante
a \(E(J)\). Una arista \(ab\) de color \(x\) da el triángulo \(abx\).
Las listas aseguran que ambas aristas cruzadas existen y están libres;
que el coloreo sea propio asegura que dos de los nuevos triángulos no
utilizan una misma arista cruzada. Los pares ficticios del núcleo no se
convierten en bloques. Una vez cubiertas todas las aristas exteriores y
del núcleo, las aristas cruzadas restantes se pueden tomar
individualmente.

\subsubsection{2.4. Otros tamaños del
núcleo}\label{otros-tamauxf1os-del-nuxfacleo}

El argumento de levantamiento no exige \(p=\lfloor(n+1)/3\rfloor\). Para
un grafo split de referencia arbitrario, con tamaños de las partes
\(p,q\), sea \(h\) el mismo parámetro local de antes.

\Needspace{13\baselineskip}
\textbf{Proposición 2.1.} Sea \(t\ge1\) un entero. Si \[
p+(4t+2)h\le q,\qquad q+t(t+5)h\le tp,
\] entonces \[
\operatorname{cp}(G)=\operatorname{cp}_{\le4}(G)
=s-u-2m=pq-\binom p2+C-A-2m.
\]

\emph{Demostración.} Se utiliza la construcción de la Sección 3 con tope
de carga \(r=th\). Para \(h\ge1\), la selección de candidatos y de pares
es válida siempre que \[
r(5h+r-2)+2(m-1)\le rp.
\] La segunda desigualdad de tamaños y \(2m\le qh\) implican esta cota
de capacidad. La primera desigualdad de tamaños es exactamente la
condición de listas \(p+2(h+2r)\le q\). Así, la misma construcción
física produce solo bloques de holgura cero y alcanza la cota inferior.
Si \(h=0\), no hay levantamiento exterior, y el coloreo con listas
completa el núcleo completo puesto que \(q\ge p\). \(\square\)

Por ejemplo, \(t=4\) permite \(p+18h\le q\le4p-36h\). Esto también cubre
el segundo tamaño extremal completo-split. Para \(n=3k+1\), ambos grafos
\[
K_k\vee I_{2k+1},\qquad K_{k+1}\vee I_{2k}
\] tienen valor de referencia \(M(n)\). Para el segundo núcleo, de
tamaño \(p=k+1\), la condición suficiente \(p\ge18h+2\) da la misma
fórmula \(M(n)+C-A-2m\). En efecto, \(q=2p-2\) satisface el intervalo
mostrado, y \(p(2p-2)-\binom p2=M(3p-2)\). La afirmación de igualdad
local de la Sección 6 siempre se refiere al grafo de referencia que en
ella se especifica.

\subsection{3. Demostración del Teorema
A}\label{demostraciuxf3n-del-teorema-a}

Si \(h=0\), entonces \(G=S\). Un coloreo propio de las aristas del
núcleo con a lo sumo \(p\) colores exteriores proporciona la
construcción triangular requerida. Un coloreo de este tipo etiqueta el
núcleo mediante \(\mathbb Z/p\mathbb Z\) y colorea \(ab\) con \(a+b\);
las aristas incidentes tienen colores distintos. Para \(p=1\) no hay
aristas del núcleo. La partición resultante tiene \(pq-\binom p2=M(n)\)
bloques y alcanza (11).

Supongamos \(h\ge1\). Se procesan las aristas exteriores \(xy\in E(H)\)
en cualquier orden, asignándoles un par disponible \(ab\) del núcleo
para obtener un \(K_4\) sobre \(\{a,b,x,y\}\). Se mantiene que los pares
del núcleo no se reutilizan, que los pares del núcleo asignados a
aristas exteriores incidentes son disjuntos en vértices y que cada
vértice del núcleo pertenece a lo sumo a \(r=4h\) bloques asignados.
Llamamos carga de un vértice al número de bloques asignados que lo
contienen.

Antes de procesar \(xy\), se excluyen los vértices del núcleo no
adyacentes a \(x\) o \(y\), los vértices utilizados antes en alguno de
los dos extremos exteriores y los vértices de carga \(r\). Las dos
primeras causas excluyen a lo sumo \[
\beta_x+\beta_y+2(d_H(x)-1)+2(d_H(y)-1)\le4h-4. \tag{14}
\] La carga total anterior a este paso es a lo sumo \(2(m-1)\), por lo
que están saturados menos de \(2m/(4h)\le q/4\) vértices. El conjunto de
candidatos tiene tamaño \(k\), con \[
\begin{aligned}
k&>p-4h+4-q/4\\
 &\ge p/2-4h+15/4\\
 &\ge5h+17/4,
\end{aligned} \tag{15}
\] utilizando \(q\le2p+1\) y (4).

Aquí un vértice saturado tiene carga exactamente \(r\). Cada arista
exterior asignada anteriormente aporta dos unidades a la carga total del
núcleo, lo que acota el número de vértices saturados. Las demás
exclusiones se estiman mediante una cota de la unión; no tienen que ser
disjuntas. Las desigualdades conjuntas de (3), en lugar de cotas
separadas de peor caso para las aristas ausentes y los grados
exteriores, dan (14).

El grafo de pares prohibidos está formado por las aristas ausentes del
núcleo y los pares del núcleo ya utilizados. Su grado máximo es a lo
sumo \(h+r=5h\). Por la Sección 2.2 y (15), existe entonces un par
disponible \(ab\). Las cuatro aristas cruzadas de \(abxy\) están
presentes y libres gracias a las exclusiones, y \(xy\) es la arista
exterior actual. Como \(a,b\) no estaban saturados, sus nuevas cargas
siguen siendo a lo sumo \(r\). Esto demuestra el paso inductivo. Si
\(m=0\), esta etapa es vacía.

Al terminar, un vértice \(a\) del núcleo de carga \(t_a\le4h\) ha
utilizado exactamente \(2t_a\) aristas cruzadas distintas. Por tanto,
\(|F_a|=\alpha_a+2t_a\le9h\), y cada arista residual del núcleo tiene
una lista de tamaño \[
|L_{ab}|\ge q-18h\ge2p-1-18h\ge p. \tag{16}
\] La Sección 2.3 completa las \(u-m\) aristas residuales del núcleo con
triángulos. Los \(m\) bloques \(K_4\) utilizan \(4m\) aristas cruzadas,
y los triángulos utilizan \(2(u-m)\). Quedan \(s-2u-2m\) aristas
cruzadas individuales. El número total de bloques es \(s-u-2m\), que
alcanza (11). Los recuentos son no negativos porque cuentan los bloques
de esta partición.

El invariante de las aristas incidentes también hace exacto el cálculo
de \(F_a\): dos bloques del levantamiento que pasan por un mismo vértice
del núcleo no pueden utilizar un mismo extremo exterior. Así, las
\(2t_a\) aristas cruzadas contadas allí son distintas, y el posterior
coloreo con listas no puede reutilizar ninguna de ellas.

Cuando \(q=2p\), la estimación de candidatos tiene más margen y la
última desigualdad de (16) solo exige \(p\ge18h\). Esto demuestra la
variante equilibrada. \(\square\)

\subsection{4. Demostración del Teorema
B}\label{demostraciuxf3n-del-teorema-b}

Todo triángulo tiene al menos una arista interna a una de las dos
partes. Hay \(u+m\) aristas de este tipo, por lo que un empaquetamiento
de triángulos tiene tamaño a lo sumo \(u+m\).

Para ver la inyección que sustenta esta cota, se elige una arista
interna de cada triángulo. Dos triángulos distintos del empaquetamiento
no pueden elegir la misma arista, porque el empaquetamiento es disjunto
en aristas. El argumento fraccional posterior es una versión ponderada
separada de esta cuenta; no supone que un empaquetamiento fraccional sea
entero. Una partición con \(t\) triángulos y todas las aristas restantes
individuales tiene \(e(G)-2t\) bloques. Por tanto, \[
\operatorname{cp}_{\le3}(G)\ge(s+u+m)-2(u+m)=s-u-m. \tag{17}
\]

Construimos \(u+m\) triángulos. El caso \(h=0\) se trató en la Sección
3. Para \(h\ge1\), se procesa primero cada arista exterior \(xy\)
asignándole una raíz \(a\) del núcleo para formar \(axy\). Se permite
que cada raíz del núcleo reciba a lo sumo \(2h\) aristas exteriores. Las
adyacencias ausentes y las asignaciones anteriores en \(x,y\) excluyen a
lo sumo \[
(\beta_x+d_H(x)-1)+(\beta_y+d_H(y)-1)\le2h-2
\] raíces. Están saturadas menos de \(m/(2h)\le q/4\) raíces. Quedan más
de \[
p-(2h-2)-q/4\ge p/2-2h+7/4>0
\] candidatos, por lo que la asignación se puede completar.

Esta etapa cubre todas las aristas exteriores y no utiliza ninguna
arista del núcleo. Cada vértice del núcleo ha utilizado a lo sumo \(4h\)
aristas cruzadas. Incluyendo sus adyacencias ausentes, prohíbe a lo sumo
\(5h\) colores exteriores. Así, cada lista de una arista del núcleo
tiene tamaño al menos \(q-10h\ge p\). La Sección 2.3 proporciona un
triángulo por cada una de las \(u\) aristas del núcleo, con lo que se
alcanzan \(u+m\) triángulos y la igualdad en (17). Junto con A, esto
demuestra (7). \(\square\)

La selección de raíces debe utilizar la información sobre las
adyacencias ausentes. Un coloreo ordinario de \(H\) seguido de una
asignación arbitraria de colores a los vértices del núcleo no aseguraría
la existencia de todas las aristas cruzadas requeridas.

La misma cuenta de aristas internas es un certificado fraccional: se
asigna peso uno a cada arista interna y cero a cada arista cruzada. Todo
triángulo tiene peso total de cobertura al menos uno. El empaquetamiento
construido alcanza ese valor, de modo que \[
\nu_3(G)=\nu_3^*(G)=u+m. \tag{18}
\]

\subsection{5. Demostración del Teorema
C}\label{demostraciuxf3n-del-teorema-c}

Cuando \(h=0\), la prescripción es vacía y se aplica la construcción
base. Supongamos \(h\ge1\) y fijemos \(\sigma\). Un \(K_4\) prescrito
necesita un par del núcleo, mientras que un \(K_5\) prescrito necesita
un triángulo del núcleo. Como antes, se exige que las aristas del núcleo
utilizadas en bloques distintos sean distintas y que los conjuntos del
núcleo asignados a aristas exteriores incidentes sean disjuntos en
vértices. Se fija como carga máxima por vértice \(r=7h\).

Antes de procesar \(xy\), las aristas cruzadas ausentes y las
asignaciones anteriores en sus extremos excluyen a lo sumo \[
\beta_x+\beta_y+3(d_H(x)-1)+3(d_H(y)-1)\le6h-6
\] vértices. La carga total anterior es a lo sumo \(3(m-1)\), por lo que
están saturados menos de \(3m/(7h)\le3q/14\) vértices. El número de
candidatos es mayor que \[
\begin{aligned}
p-6h+6-3q/14
&\ge4p/7-6h+81/14\\
&\ge30h+81/14>30h+3.
\end{aligned} \tag{19}
\] Cada participación anterior utilizó a lo sumo dos aristas incidentes
del núcleo. Incluyendo las aristas ausentes del núcleo, el grafo de
pares prohibidos tiene grado máximo a lo sumo \(h+2r=15h\). La Sección
2.2 proporciona un triángulo disponible del núcleo. Se utiliza entero
para la etiqueta 5 y solo uno de sus pares para la etiqueta 4, cargando
únicamente los vértices que se usan. Las mismas exclusiones de la
Sección 3 aseguran que todas las aristas del bloque estén presentes,
libres y sean compatibles con los invariantes.

Una terna se utiliza únicamente para encontrar una configuración
disponible. Para un \(K_4\) prescrito, su tercer vértice no utilizado no
se añade al levantamiento ni recibe carga. Para un \(K_5\), una
participación utiliza dos aristas incidentes del núcleo; por eso el
grado de pares prohibidos se acota por \(h+2r\) en el paso siguiente.

Al terminar, cada vértice del núcleo ha utilizado a lo sumo \(2r=14h\)
aristas cruzadas. Las listas de completación tienen tamaño al menos \[
q-2h-4r=q-30h\ge2p-1-30h\ge p,
\] puesto que \(p\ge63h\) y \(h\ge1\). Se aplica la Sección 2.3 y se
dejan las aristas cruzadas restantes como bloques de una sola arista.
Todos los bloques tienen holgura cero en (11), lo que demuestra la
optimalidad.

Si \(k\) aristas exteriores tienen etiqueta 5, los bloques grandes
utilizan \(m+2k\) aristas del núcleo y \(4m+2k\) aristas cruzadas. Así,
\(N_3=u-m-2k\) y \[
N_2=s-(4m+2k)-2(u-m-2k)=s-2u-2m+2k,
\] lo que demuestra (9). Por último, el orden del único bloque que
contiene cada arista exterior permite recuperar su etiqueta. Por tanto,
prescripciones distintas producen particiones distintas, lo que da al
menos \(2^m\) óptimos. \(\square\)

\subsection{6. Estabilidad local para grafos
cordales}\label{estabilidad-local-para-grafos-cordales}

\textbf{Lema 6.1 (rigidez del núcleo).} Bajo las hipótesis de A, si
\(G\) es cordal, entonces \(C=0\).

\emph{Demostración.} El caso \(h=0\) es inmediato. En otro caso,
supongamos que falta una arista \(ab\) del núcleo. Sus extremos tienen
al menos \[
q-\alpha_a-\alpha_b\ge q-2h\ge2p-1-2h>h+1
\] vecinos exteriores comunes. Como \(\Delta(H)\le h\), este conjunto no
es un clique. Se eligen vecinos comunes no adyacentes \(x,y\). Los
vértices \(a,x,b,y\) inducen un ciclo de longitud cuatro: sus dos
posibles cuerdas \(ab,xy\) están ausentes. Esto contradice la
cordalidad. \(\square\)

En consecuencia, para grafos cordales de este entorno, A da \[
\operatorname{cp}(G)=M(n)-A-2m. \tag{20}
\] La distancia de edición etiquetada es
\(d_{\rm edit}=|E(G)\triangle E(S)|=A+m\). Por tanto, \[
d_{\rm edit}\le M(n)-\operatorname{cp}(G)\le2d_{\rm edit}. \tag{21}
\] En particular, la igualdad \(\operatorname{cp}(G)=M(n)\) obliga a
\(G=S\) dentro de este entorno. Esta es una conclusión local relativa al
grafo de referencia fijo; no caracteriza todos los grafos cordales
extremales.

\subsection{7. Optimalidad fraccional y tipos de bloques
óptimos}\label{optimalidad-fraccional-y-tipos-de-bloques-uxf3ptimos}

Definamos el número de partición fraccional
\(\operatorname{cp}_{=}^{*}(G)\) minimizando \(\sum_Q x_Q\), sobre
cliques de orden al menos dos, sujeto a \(x_Q\ge0\) y
\(\sum_{Q\ni e}x_Q=1\) para toda arista presente \(e\). Asignemos a las
aristas presentes cruzadas, del núcleo y exteriores los pesos
respectivos \(1,-1,-2\). Su suma en un clique de tipo \((a,b)\) es
\(1-\ell(a,b)\le1\), y su suma total es \(s-u-2m\). Así, toda partición
fraccional factible satisface \[
\sum_Qx_Q-(s-u-2m)=\sum_Q\ell(a_Q,b_Q)x_Q\ge0.
\] Bajo A, la construcción entera alcanza la cota, lo que demuestra \[
\operatorname{cp}_{=}^{*}(G)=\operatorname{cp}_{\le4,=}^{*}(G)
=\operatorname{cp}_{\le4}(G)=\operatorname{cp}(G). \tag{22}
\] Aquí el subíndice \(\le4\) restringe también las variables de los
bloques fraccionales. Esta es una igualdad para el objetivo uniforme de
partición; no afirma que todo el politopo sea entero. Las restricciones
de aristas son igualdades, de modo que este problema también es distinto
del de cobertura fraccional.

De manera equivalente, se maximiza la ganancia \(\binom{|Q|}{2}-1\)
sobre empaquetamientos de cliques disjuntos en aristas, o sobre
empaquetamientos fraccionales con capacidades unitarias por arista.
Completar la capacidad no utilizada mediante \(K_2\) muestra que ambos
óptimos de ganancia son \(2u+3m\). Esto proporciona un caso exacto de la
distinción del Artículo V entre un certificado de empaquetamiento y su
realización. La construcción y la cota inferior de este trabajo
demuestran directamente la igualdad; no se necesita ningún teorema de
transferencia del Artículo V.

La ecuación (11) también muestra que toda partición entera óptima
utiliza únicamente los cuatro tipos de (12). La misma conclusión vale
para el soporte de toda partición fraccional óptima. Ningún bloque
óptimo tiene orden mayor que cinco, incluso si el grafo contiene cliques
mayores. El Teorema A proporciona un óptimo con orden a lo sumo cuatro,
mientras que C permite las elecciones independientes indicadas entre
cuatro y cinco. El Apéndice B cuantifica cuántas aristas pueden
pertenecer a otros tipos en una partición casi óptima.

\subsection{8. Reconocer el núcleo y permitir un vértice
excepcional}\label{reconocer-el-nuxfacleo-y-permitir-un-vuxe9rtice-excepcional}

La fórmula del Teorema A se enuncia respecto de un núcleo de referencia.
Primero demostramos que, en esa región local, el núcleo se puede
recuperar a partir del grafo. Después examinamos una familia con un
vértice excepcional. En esa familia, el coste adicional tiene una
fórmula exacta, y cambiar el núcleo no necesariamente lo elimina.

\subsubsection{8.1. Una fórmula intrínseca en la región
local}\label{una-fuxf3rmula-intruxednseca-en-la-regiuxf3n-local}

Para todo conjunto de vértices \(Q\), sin condición sobre su tamaño,
pongamos \[
\Phi_Q(G)=e_G(Q,V\setminus Q)-e(G[Q])-2e(G[V\setminus Q]).
\] La identidad de holgura de la Sección 2 da
\(\Phi_Q(G)\le\operatorname{cp}(G)\) para todo \(Q\). Bajo las hipótesis
del Teorema A, se alcanza la igualdad en el núcleo de referencia. Por
tanto, \[
\operatorname{cp}(G)=\max_{Q\subseteq V}\Phi_Q(G)=\Phi_P(G).
\] Esta conclusión no afirma que el conjunto maximizador sea único.

Si \(n\ge3\), las mismas hipótesis dan la identificación más útil \[
P=\{v\in V:2d_G(v)>n\}.
\] En efecto, un vértice del núcleo tiene grado al menos \(n-1-h\),
mientras que un vértice exterior tiene grado a lo sumo \(p+h\). Para
\(h=0\), la separación afirmada se sigue de \(n\ge3\) y del valor
especificado de \(p\). Para \(h\ge1\), se utilizan \(p\ge18h+1\) y
\(n\ge3p-1\) para obtener \(n-1-h>n/2>p+h\).

Así, el valor se puede calcular sin recibir la partición de referencia:
se forma el conjunto de vértices por encima de este umbral de grado, se
comprueba su tamaño y su grado de modificación local y, si se aplica el
Teorema A, se cuentan sus tres clases de aristas. Con una matriz de
adyacencia, estas operaciones requieren tiempo \(O(n^2)\). Este es un
algoritmo para el valor en la región reconocida. No proporciona una cota
del tiempo necesario para construir la partición mediante el teorema
exacto de coloreo con listas.

\subsubsection{8.2. El defecto exacto de un vértice
excepcional}\label{el-defecto-exacto-de-un-vuxe9rtice-excepcional}

En esta subsección se permiten tamaños arbitrarios \(p,q\). Sea \(P\) un
clique de tamaño \(p\), sea \(R\) un conjunto independiente de tamaño
\(q-1\) e incluyamos todas las aristas entre \(P\) y \(R\). Se añade un
vértice \(z\) con exactamente \(a\) vecinos en \(P\) y \(b\) vecinos en
\(R\). Supongamos \(q\ge p+3\); se permiten valores cero de \(p,a,b\)
cuando sean compatibles con estas condiciones. Respecto del núcleo
\(P\), escribamos \[
u=\binom p2,\qquad m=b,\qquad s=p(q-1)+a,\qquad \Phi=s-u-2b,
\] y definamos \[
D(a,b)=\max\{0,2b-a,3b-2a\}.
\] La subfamilia \(a=p\) es split. El caso general no tiene por qué ser
cordal: si \(a<p\) y \(b\ge2\), un no vecino de \(z\) en \(P\) y dos
vecinos de \(z\) en \(R\) forman un ciclo inducido de longitud cuatro
con \(z\).

\textbf{Teorema D.} Para este grafo, \[
\operatorname{cp}(G)=\operatorname{cp}_{\le4}(G)
=\operatorname{cp}^{*}_{=}(G)=\Phi+D(a,b).
\] Además, \[
\nu_3(G)=\nu_3^*(G)=u+\min(a,b),\qquad
\operatorname{cp}_{\le3}(G)-\operatorname{cp}(G)
=\min\{b,\max(0,a-b)\}.
\] Aquí \(\operatorname{cp}^{*}_{=}\) denota partición fraccional con
igualdad en cada arista, como en la Sección 7, y \(\nu_3\) cuenta
triángulos disjuntos en aristas.

\emph{Demostración.} Todo bloque que contiene una arista de \(z\) a
\(R\) contiene exactamente un vértice de \(R\), el vértice \(z\) y un
número \(k\) de vecinos de \(z\) en \(P\). Hay exactamente \(b\) bloques
de este tipo, puesto que \(R\) es independiente. Sus conjuntos del
núcleo son disjuntos dos a dos: compartir un vértice del núcleo
reutilizaría su arista hacia \(z\). Por tanto, la suma de sus tamaños en
el núcleo es a lo sumo \(a\).

La holgura de uno de estos bloques es \[
\ell(k,2)=3-2k+\binom k2\ge\max\{0,2-k,3-2k\}.
\] Todos los demás bloques tienen holgura no negativa. Al sumar se
obtiene \(\operatorname{cp}(G)\ge\Phi+D(a,b)\). La misma cota es válida
en el caso fraccional. Para verlo directamente, se parte de los pesos de
aristas de la Sección 7. Además de esos pesos, se utiliza cada una de
las dos correcciones \[
(-1,2),\qquad (-2,3)
\] en las aristas \(zP,zR\), respectivamente. Cada peso corregido es
factible para todo clique: las desigualdades mostradas tratan los
bloques que intersectan \(R\), y la corrección es no positiva en los
bloques que pasan por \(z\) y que, aparte de ese vértice, solo
intersectan \(P\). Los tres valores duales son
\(\Phi,\Phi+2b-a,\Phi+3b-2a\).

Este argumento dual es necesario para la conclusión fraccional. Contar
los \(b\) bloques reales que pasan por \(zR\) demuestra la cota inferior
entera, pero esa cuenta por sí sola no justificaría un enunciado sobre
pesos fraccionales arbitrarios. La factibilidad de las tres asignaciones
de pesos a las aristas proporciona la cota inferior fraccional de manera
independiente de la construcción entera.

Para la construcción que alcanza la cota, se asignan subconjuntos
disjuntos de los \(a\) vecinos de \(z\) en el núcleo a sus \(b\) vecinos
en \(R\). Si \(a\le b\), se utilizan \(a\) triángulos y \(b-a\) aristas
individuales que pasan por \(zR\). Si \(b\le a\le2b\), se utilizan
\(a-b\) copias de \(K_4\) y \(2b-a\) triángulos. Si \(a\ge2b\), se
utilizan \(b\) copias de \(K_4\) y se cubren individualmente las aristas
\(zP\) no utilizadas. Estas elecciones cubren todas las aristas
incidentes con \(z\), y su holgura total es \(D(a,b)\).

Cada vértice del núcleo ha utilizado a lo sumo una arista cruzada hacia
\(R\). Se excluye \(z\) de los colores de completación. Cada arista
residual del núcleo tiene entonces al menos \(|R|-2=q-3\ge p\) colores
disponibles. La completación mediante coloreo con listas de la Sección 2
cubre el núcleo residual con triángulos; todas las aristas cruzadas
restantes son bloques individuales. Esos bloques finales tienen holgura
cero. Así se construye un óptimo de orden a lo sumo cuatro, y su función
indicadora también alcanza la cota fraccional.

Para los triángulos, se asigna peso de cobertura uno a cada arista del
núcleo y, según cuál de los dos conjuntos sea menor, a todas las aristas
\(zP\) o a todas las aristas \(zR\). Todo triángulo recibe peso al menos
uno. El peso de cobertura es \(u+\min(a,b)\), lo que demuestra la cota
superior fraccional. Se utilizan \(\min(a,b)\) triángulos que pasan por
\(z\), con vecinos distintos en el núcleo y en el exterior. Se completan
todas las \(u\) aristas del núcleo con colores en \(R\), utilizando de
nuevo listas de tamaño al menos \(p\). Esto alcanza el peso de
cobertura. Por último, un empaquetamiento de \(t\) triángulos se
completa a \(e(G)-2t\) piezas. Al sustituir las dos fórmulas exactas en
los tres intervalos de \(a\), se obtiene la brecha de restricción
enunciada. \(\square\)

\subsubsection{8.3. Lo que cambiar el núcleo no puede
reparar}\label{lo-que-cambiar-el-nuxfacleo-no-puede-reparar}

La familia excepcional ya impide intervalos locales mayores. Tomemos
\(a=p\) y \(b=\lfloor p/2\rfloor+1\), con \(p\ge1\) y la hipótesis
vigente \(q\ge p+3\) de la Sección 8.2. El grafo es split y, por tanto,
cordal, y su grado de modificación respecto del núcleo original es
\(h=b\). Su defecto es \(D(p,b)=2b-p>0\). Así, una fórmula sin defecto
para toda perturbación cordal no puede extenderse a razones de
modificación estrictamente mayores que \(1/2\). Esto proporciona una
obstrucción superior; no se demuestra un umbral suficiente coincidente.
Para la elección \(q=2p\), que exige \(p\ge3\) bajo esta hipótesis, el
vértice excepcional satisface \(d_G(z)=p+b>n/2\). El criterio de grados
de la Sección 8.1 lo situaría en el núcleo, pero sus hipótesis locales
fallan: respecto de \(P\cup\{z\}\), el vértice \(z\) tiene \(2p-1-b\)
aristas cruzadas ausentes.

Se podría intentar eliminar el defecto cambiando el núcleo de
referencia. Incluso optimizar sobre todas las biparticiones puede no
bastar. Si \(q=2p\) y \(p\ge19\), entonces \[
\max_{Q\subseteq V}\Phi_Q(G)=\Phi+\max\{0,3b-2a\}.
\] He aquí una comparación directa. Para un conjunto propuesto \(Q\),
pongamos \(x=|Q\cap P|\), \(x'=p-x\), \(y=|Q\cap R|\) y \(y'=2p-1-y\).
Los vértices normales aportan al menos \[
2\left(\binom x2+xy\right)
+3\left(\binom{x'}2+x'y'\right)
\ge 2\binom p2+3x'+3y
\] a \(2e(G[Q])+3e(G[V\setminus Q])\). Para comprobar la desigualdad, el
doble de su diferencia es \(E=5x'^2+(8p-13)x'+(4p-6)y-10x'y\), y \[
\begin{aligned}
(2p-1)E&=(2p-1-y)\bigl(5x'^2+(8p-13)x'\bigr)\\
&\quad+y\bigl(p^2+2px+5x^2-19p+3x+6\bigr)\ge0.
\end{aligned}
\] Aquí se utiliza la condición \(p\ge19\) mediante \(p^2-19p\ge0\); es
una condición suficiente para esta estimación. Si \(z\in Q\), al menos
\(a-x'\) de sus vecinos del núcleo están en \(Q\), lo que da un coste
total al menos \(2\binom p2+2a\). Si \(z\notin Q\), al menos \(b-y\) de
sus vecinos en \(R\) están fuera de \(Q\), lo que da un coste al menos
\(2\binom p2+3b\). Como \(\Phi_Q=e(G)-2e(G[Q])-3e(G[V\setminus Q])\),
estas cotas dan el máximo afirmado. Lo alcanzan \(P\) o \(P\cup\{z\}\),
cuyos valores difieren en \(3b-2a\).

Para \(a=p\) y \(0\le b\le p\), la brecha es \[
\operatorname{cp}(G)-\max_Q\Phi_Q(G)
=\max\{0,2b-p\}-\max\{0,3b-2p\}.
\] Ya es positiva cuando \(p/2<b<2p/3\). El valor \(2p/3\) marca el
cambio entre los dos núcleos de referencia maximizadores y el máximo de
la expresión lineal por tramos de la brecha, y no el comienzo del fallo
después de optimizar el núcleo. En particular, para \(a=p\), \(3b=2p\) y
\(p\ge19\), \[
\operatorname{cp}(G)-\max_Q\Phi_Q(G)=p/3.
\] Todas las modificaciones son incidentes con \(z\); todos los demás
vértices tienen grado de modificación a lo sumo uno. Tomando \(p=3t\),
\(t\ge7\) se obtienen ejemplos cordales concretos con brecha \(t\). Por
tanto, un número acotado de vértices excepcionales no se puede tratar
como un coste aditivo constante sin otra hipótesis. Los tres duales del
Teorema D detectan una obstrucción de recursos que la familia de
biparticiones por sí sola no ve. En la terminología del Artículo V, las
demandas exteriores que compiten comparten el conjunto finito de aristas
utilizables desde \(z\) hasta el núcleo; una distancia total de edición
pequeña no acota la pérdida de realización resultante.

\subsection{9. Prescripciones y número de particiones
óptimas}\label{prescripciones-y-nuxfamero-de-particiones-uxf3ptimas}

Sea \(N(G)\) el número de particiones en cliques del grafo etiquetado
\(G\) que alcanzan \(\operatorname{cp}(G)\). Las particiones son
familias de conjuntos de vértices, y no listas ordenadas de bloques ni
clases de isomorfismo.

\textbf{Teorema E.} Bajo las hipótesis del Teorema C, si \(p\ge16\),
cada prescripción fija de aristas exteriores se realiza mediante al
menos \(p^{\lfloor p^2/64\rfloor}\) particiones óptimas. En
consecuencia, \[
2^m p^{\lfloor p^2/64\rfloor}\le N(G)\le(n+1)^{5e(G)}.
\] En particular, \(\log N(G)=\Theta(p^2\log p)\) en esta región cuando
\(p\to\infty\).

La demostración, presentada en el Apéndice C, fija el levantamiento
exterior y cuenta muchos coloreos compatibles de un subgrafo ralo
suficientemente grande del núcleo residual. El teorema de Vizing extrae
ese subgrafo; Häggkvist--Janssen extiende cada uno de los prefijos
contados. Estos dos resultados de coloreo tienen funciones distintas. El
núcleo residual tiene grado cercano a \(p\), por lo que tratarlo como un
grafo de grado \(O(h)\) no justificaría la conclusión.

El levantamiento también permite un cambio acotado cuando se prescribe
un bloque exterior físico. Damos el invariante y el argumento, puesto
que esta conclusión concierne a algo más que elegir los órdenes de los
bloques.

Para una arista exterior \(e=xy\), un \emph{conjunto legal del núcleo}
es un subconjunto \(N\subseteq P\) que induce un clique y cuyos vértices
son todos adyacentes tanto a \(x\) como a \(y\). La legalidad se refiere
aquí a ese bloque propuesto; la compatibilidad con el levantamiento
actual se impone o se repara por separado.

\textbf{Proposición 9.1 (extensión y cambio del levantamiento
exterior).} Supongamos las hipótesis del Teorema C. Para cada arista
exterior asignada \(e=xy\), sea \(T_e\subseteq P\) un conjunto de tamaño
dos o tres que induce un clique y es adyacente a ambos extremos de
\(e\). Se exige que los conjuntos de aristas del núcleo de los \(T_e\)
sean disjuntos dos a dos, que \(T_e\cap T_f=\varnothing\) siempre que
\(e,f\) sean aristas exteriores distintas e incidentes y que cada
vértice del núcleo pertenezca a lo sumo a \(7h\) conjuntos asignados.
Todo levantamiento parcial de este tipo se extiende a una partición
óptima con los órdenes especificados en todas las aristas exteriores,
siempre que esos órdenes concuerden con los conjuntos asignados. Dado un
levantamiento completo y un nuevo conjunto legal \(N\) del núcleo de
tamaño dos o tres para una arista \(e_0\), existe un levantamiento
completo que contiene \(N\) en \(e_0\), que cambia a lo sumo trece
asignaciones, incluida la de \(e_0\), y conserva la cota de carga. Solo
puede cambiar el orden prescrito en \(e_0\).

\emph{Demostración.} Las estimaciones de la Sección 5 dependen del
número de aristas asignadas y del invariante indicado, por lo que su
inducción puede comenzar desde el levantamiento parcial prescrito. La
completación por coloreo con listas conserva entonces todos los bloques
asignados.

Para efectuar el cambio, se elimina la asignación antigua en \(e_0\).
También se elimina toda asignación que utilice un par de \(N\) o un
vértice de \(N\) en un extremo de \(e_0\). Cada uno de los a lo sumo
tres pares de \(N\) tiene a lo sumo un propietario, porque las aristas
del núcleo son disjuntas. Cada una de las a lo sumo seis incidencias
extremo--núcleo tiene a lo sumo un propietario, por la condición de las
aristas incidentes. Después de estas eliminaciones, si un vértice de
\(N\) aún tiene carga \(7h\), se elimina otra asignación que lo
contenga. Bastan a lo sumo tres eliminaciones adicionales. Por tanto, se
han eliminado a lo sumo \(1+3+6+3=13\) asignaciones exteriores. Como
\(e_0\) es una arista exterior, \(h\ge1\), de modo que un vértice
saturado tiene una asignación disponible para eliminar. Se inserta el
nuevo bloque: no tiene conflictos de pares ni de aristas cruzadas, y sus
vértices del núcleo tenían carga menor que \(7h\). Se extiende el
levantamiento parcial resultante mediante la primera afirmación. Se
conservan todas las asignaciones que quedan fuera del conjunto
eliminado. \(\square\)

La completación triangular puede cambiar globalmente. La Proposición 9.1
acota los cambios en la capa de levantamiento exterior; no implica un
cambio de trece bloques en la partición completa ni la conectividad de
particiones óptimas arbitrarias. La subclase de joins completos con
levantamientos formados solo por \(K_5\) admite las conclusiones más
fuertes de la Sección 10.

El papel de la serie queda ahora más explícito. El Artículo IV {[}IV{]}
utiliza localización estructural y cuentas terminales físicas para
cubrir el presupuesto global de partición. Los resultados de este
trabajo determinan óptimos exactos y el coste de una familia terminal
restringida en una región local, mientras que el Teorema D identifica
una obstrucción adicional cuando falla el control local. El Artículo V
{[}V{]} proporciona un lenguaje para contabilizar las pérdidas de
realización, pero su desigualdad abstracta de transferencia no construye
la raíz ni las realizaciones de bajo coste utilizadas aquí. Ninguna de
estas conclusiones locales se infiere únicamente de una distancia total
de edición pequeña.

\subsection{10. Completación estable y conectividad restringida en joins
completos}\label{completaciuxf3n-estable-y-conectividad-restringida-en-joins-completos}

El cambio exterior de la Proposición 9.1 controla las asignaciones,
mientras que una partición en cliques también contiene triángulos
terminales y aristas cruzadas individuales. En un join completo con
todos los bloques exteriores de orden cinco, estos bloques restantes se
pueden reparar con un coste proporcional al número de asignaciones
cambiadas. La razón es una compensación exacta entre las aristas del
núcleo y las aristas cruzadas utilizadas por un levantamiento.

En toda esta sección, sea \(p\ge1\), sea \(G=K_P\vee H\), donde
\(|P|=p\), \(|V(H)|=q\), y escribamos \(u=\binom p2\), \(m=e(H)\). Un
\emph{levantamiento completo legal formado solo por \(K_5\)} asigna a
cada arista exterior \(xy\) una terna \(T_{xy}\) del núcleo, de modo que
los bloques resultantes \(T_{xy}\cup\{x,y\}\) son disjuntos en aristas
dos a dos. Su carga en un vértice del núcleo es el número de estos
bloques que contienen ese vértice. Un triángulo terminal tiene tipo
\((2,1)\).

\subsubsection{10.1. Completación a partir de un prefijo legal
arbitrario}\label{completaciuxf3n-a-partir-de-un-prefijo-legal-arbitrario}

\textbf{Lema 10.1.} Supongamos \(q\ge2p-1\) y fijemos un levantamiento
completo legal formado solo por \(K_5\). Toda familia de triángulos
terminales disjuntos en aristas que no comparte aristas con sus bloques
se extiende a una partición óptima en cliques que contiene el
levantamiento. Los bloques restantes son triángulos terminales y aristas
cruzadas individuales.

\emph{Demostración.} Sea \(J\) el grafo del núcleo no utilizado y sea
\(F_a\) el conjunto de vértices exteriores cuyas aristas cruzadas en
\(a\) utiliza el levantamiento. Si \(a\) tiene carga \(t_a\), la
disjunción en aristas da \[
d_J(a)=p-1-2t_a,\qquad |F_a|=2t_a,
\qquad d_J(a)+|F_a|=p-1.
\] Los triángulos terminales prescritos determinan un coloreo propio
parcial de las aristas de \(J\), con los vértices exteriores como
colores. El color de \(ab\) evita \(F_a\cup F_b\). Para una arista no
coloreada \(ab\), el levantamiento y los colores ya utilizados en sus
extremos prohíben a lo sumo \[
|F_a|+|F_b|+(d_J(a)-1)+(d_J(b)-1)=2p-4<q
\] colores. Así, toda arista restante se puede colorear de forma voraz
sin cambiar el prefijo. Si existe una arista no coloreada, entonces
\(p\ge2\) y ambos extremos tienen grado positivo, de modo que la cuenta
no contiene ningún término de grado negativo. Se convierten los colores
resultantes en triángulos terminales y se toman individualmente las
aristas cruzadas restantes. Todos los bloques tienen holgura cero en
(11), lo que demuestra la optimalidad. En particular, los números de
bloques son \(m\), \(u-3m\) y \(pq-2u\), respectivamente. \(\square\)

El lema supone que el levantamiento completo ya existe. No requiere una
cota de su carga ni una cota superior de \(q\). Su prefijo puede
contener cualesquiera triángulos terminales compatibles, y no solo
triángulos seleccionados por una ejecución anterior del procedimiento
voraz.

\subsubsection{10.2. Reparar la partición
completa}\label{reparar-la-particiuxf3n-completa}

\textbf{Proposición 10.2.} Supongamos \(q\ge2p-1\). Sean \(L_0,L_1\)
levantamientos completos legales formados solo por \(K_5\) cuyas
asignaciones difieren en exactamente \(k\) aristas exteriores. Para toda
partición óptima en cliques \(\Pi_0\) que extiende \(L_0\), existe una
partición óptima \(\Pi_1\) que extiende \(L_1\) tal que \[
|\Pi_0\setminus\Pi_1|=|\Pi_1\setminus\Pi_0|\le34k,
\qquad |\Pi_0\triangle\Pi_1|\le68k.
\] La partición \(\Pi_1\) puede conservar todos los triángulos
terminales de \(\Pi_0\) cuyas aristas no pertenecen a ningún bloque de
\(L_1\).

\emph{Demostración.} El Lema 10.1 muestra que se alcanza la cota
inferior \(pq-u-2m\). Por tanto, la identidad de holgura obliga a que
toda partición óptima utilice únicamente los tipos de holgura cero. Como
\(\Pi_0\) contiene un bloque del levantamiento para cada arista
exterior, todos sus demás bloques son triángulos terminales o aristas
cruzadas individuales. Su familia terminal \(\mathcal T_0\) tiene
\(u-3m\) miembros.

Se conservan los miembros de \(\mathcal T_0\) que no comparten aristas
con ningún bloque de \(L_1\), y sea \(R\) el número de los descartados.
Un bloque del levantamiento que no cambia no puede compartir aristas con
un triángulo terminal antiguo. Cada bloque nuevo cambiado tiene tres
aristas del núcleo y seis aristas cruzadas que pueden pertenecer a
triángulos terminales antiguos. Como esos triángulos son disjuntos en
aristas, se descartan a lo sumo nueve por asignación cambiada. Por
tanto, \(R\le9k\). Los triángulos conservados forman un prefijo legal
para \(L_1\). Ambos grafos residuales del núcleo tienen \(u-3m\)
aristas, por lo que el Lema 10.1 añade exactamente \(R\) triángulos
terminales a este prefijo.

Los grafos residuales del núcleo no tienen por qué tener los mismos
conjuntos de aristas. Solo se comparan sus cardinalidades: cada
triángulo conservado representa una arista distinta del nuevo núcleo
residual, y cada una de las demás aristas es una nueva demanda. Esto
explica por qué el número de demandas que se deben cubrir de nuevo es
exactamente el número de triángulos descartados.

Se pierden a lo sumo \(k\) bloques antiguos del levantamiento y \(R\)
triángulos terminales antiguos. Una arista cruzada individual antigua
solo puede desaparecer al incorporarse a un bloque nuevo cambiado del
levantamiento o a un nuevo triángulo terminal. Estos utilizan a lo sumo
\(6k+2R\) aristas cruzadas. Así, \[
|\Pi_0\setminus\Pi_1|\le k+R+6k+2R
=7k+3R\le34k.
\] Las dos particiones óptimas tienen el mismo tamaño, por lo que sus
diferencias unilaterales son iguales. Su suma da la cota de diferencia
simétrica. Esto también cubre \(k=0\), cuando se conserva la partición
original. \(\square\)

Por tanto, para los joins completos de carga acotada de la Proposición
9.1, se puede insertar un nuevo \(K_5\) legal prescrito en una arista
exterior perdiendo a lo sumo \(34\cdot13=442\) bloques antiguos y
cambiando a lo sumo \(68\cdot13=884\) bloques en diferencia simétrica. A
diferencia de la Proposición 9.1 por sí sola, esto controla la partición
completa. Las hipótesis de join completo y de utilizar solo \(K_5\) se
usan en el Lema 10.1.

\subsubsection{\texorpdfstring{10.3. Caminos entre óptimos de carga
acotada con levantamientos solo de
\(K_5\)}{10.3. Caminos entre óptimos de carga acotada con levantamientos solo de K\_5}}\label{caminos-entre-uxf3ptimos-de-carga-acotada-con-levantamientos-solo-de-k_5}

\textbf{Teorema 10.3.} Supongamos \[
q\in\{2p-1,2p,2p+1\},\qquad h=\Delta(H),\qquad p\ge63h.
\] Sea \(\mathcal C\) la clase de particiones óptimas en cliques en las
que toda arista exterior pertenece a un bloque \(K_5\) de tipo \((3,2)\)
y el levantamiento inducido tiene carga a lo sumo \(7h\) en cada vértice
del núcleo. Cualesquiera dos miembros de \(\mathcal C\) están unidos
dentro de \(\mathcal C\) por una sucesión de a lo sumo \(u-2m\) pasos,
cada uno de diferencia simétrica de bloques a lo sumo \(884\). Más
precisamente, a lo sumo \(m\) pasos cambian el levantamiento exterior; a
lo sumo \(u-3m\) pasos adicionales lo mantienen fijo y tienen diferencia
simétrica a lo sumo \(18\).

\emph{Demostración.} Se fijan la partición objetivo \(\Pi_*\) y su
levantamiento \(L_*\). Primero se alinean las asignaciones exteriores.
Se protege cada asignación del levantamiento actual que ya coincide con
su objetivo. En una arista exterior no alineada \(e_0\), se prescribe la
terna objetivo \(N=L_*(e_0)\).

Como en la Proposición 9.1, se eliminan la asignación antigua en
\(e_0\), los propietarios de los pares del núcleo de \(N\) y los
propietarios de las aristas cruzadas que necesita el bloque objetivo. Se
eliminan a lo sumo \(1+3+6\) asignaciones. Ninguna está protegida: un
propietario protegido entraría en conflicto con \(L_*(e_0)\) en el
levantamiento objetivo legal. Si un vértice de \(N\) sigue saturado con
carga \(7h\), a lo sumo \(7h-1\) de sus asignaciones pueden estar
protegidas. En efecto, el levantamiento objetivo tiene carga a lo sumo
\(7h\) en ese vértice y también contiene su asignación adicional en
\(e_0\). Por tanto, hay una asignación no protegida que se puede
eliminar. A lo sumo tres eliminaciones adicionales dejan espacio para el
bloque objetivo.

Se inserta ese bloque y se extiende el levantamiento parcial conservado
mediante la Proposición 9.1. Todas las asignaciones siguen siendo de
orden cinco, se mantiene la cota de carga y cambian a lo sumo trece
asignaciones. Se conservan todas las asignaciones protegidas y \(e_0\)
queda alineada. La Proposición 10.2 proporciona una partición completa
óptima a distancia entre bloques a lo sumo \(884\). A lo sumo \(m\)
repeticiones hacen que el levantamiento sea igual a \(L_*\). Si \(m>0\),
entonces \(h\ge1\), por lo que el umbral de saturación utilizado aquí es
positivo; si \(h=0\), esta fase es vacía.

Ahora se mantiene fijo este levantamiento y se alinean los triángulos
terminales. Estos son coloreos propios con listas de las aristas de su
grafo residual del núcleo \(J\), y \(|E(J)|=u-3m\). Se protege cada
triángulo que ya se comparte con \(\Pi_*\). Se elige una arista del
núcleo no alineada \(ab\) cuyo color objetivo es \(x\). Se eliminan su
triángulo actual y los triángulos, si los hay, en las a lo sumo dos
aristas incidentes que tienen actualmente color \(x\). Se eliminan a lo
sumo \(R\le3\) triángulos terminales. Ninguno está protegido, puesto que
el coloreo objetivo es propio.

Se inserta \(abx\). No comparte aristas con el levantamiento ni con
ningún triángulo conservado, de modo que el Lema 10.1 completa este
prefijo legal manteniendo todos los triángulos protegidos. Se adquiere
al menos un nuevo triángulo objetivo. Se pierden a lo sumo \(R\)
triángulos terminales antiguos, lo que da una diferencia simétrica
terminal a lo sumo \(2R\). Solo las a lo sumo \(2R\) aristas cruzadas de
los nuevos triángulos terminales pueden consumir aristas cruzadas
individuales antiguas. Ambas particiones tienen \(pq-2u\) bloques de una
sola arista cruzada, por lo que su diferencia simétrica es a lo sumo
\(4R\). Por tanto, el paso completo cambia a lo sumo \(6R\le18\)
bloques.

Después de a lo sumo \(u-3m\) pasos, todos los triángulos terminales
coinciden con el objetivo. Las aristas cruzadas individuales restantes
quedan entonces determinadas, por lo que las particiones completas
coinciden. Al concatenar las dos fases se obtienen a lo sumo
\(m+(u-3m)=u-2m\) pasos. La desigualdad \(u\ge3m\) se sigue de los
triángulos del núcleo disjuntos en aristas de cualquier levantamiento
legal formado solo por \(K_5\). \(\square\)

La construcción del Teorema C proporciona miembros de \(\mathcal C\).
Para \(p=1\) o \(2\), un levantamiento completo formado solo por \(K_5\)
obliga a \(m=0\); la primera fase es vacía y se aplica el mismo
argumento terminal. El teorema concierne a esta clase de carga acotada
con levantamientos formados solo por \(K_5\). No proporciona la
conectividad de óptimos mixtos \(K_4/K_5\) ni una reparación tras
cambios arbitrarios del grafo. No se afirma que las constantes \(18\),
\(68\) o \(884\) sean óptimas.

\textbf{Formalización y reproducibilidad.} El complemento
{[}VI-companion{]} registra el punto de entrada canónico en Lean 4.35,
la correspondencia de teoremas, las dependencias demostradas y la
compilación final con caché y auditoría de axiomas de A--E, los
resultados intrínsecos y excepcionales y la Sección 10. Las particiones
y los empaquetamientos triangulares fraccionales usan pesos reales.
Häggkvist--Janssen y Vizing son dependencias demostradas {[}HJ-formal,
Viz-formal{]}. La auditoría final permite únicamente los axiomas
fundacionales habituales de Lean. Los ejemplos finitos y sus
certificados, comprobados independientemente, constituyen evidencia de
regresión separada de las demostraciones universales. El complemento
indica los dominios ensayados, las dependencias de las demostraciones y
el alcance de la verificación.

\subsection{Apéndice A. Una variante elemental para joins
completos}\label{apuxe9ndice-a.-una-variante-elemental-para-joins-completos}

\textbf{Teorema A.1.} Sean \(q\in\{2p-1,2p,2p+1\}\), \(n=p+q\) y
\(G=K_p\vee H\), donde \(|H|=q\) y \(d=\Delta(H)\). Si \(p\ge24d\),
entonces \[
\operatorname{cp}(K_p\vee H)=M(n)-2e(H). \tag{23}
\] Existe un óptimo en el que toda arista exterior pertenece a un
\(K_5\) de tipo \((3,2)\).

\emph{Demostración.} Para \(d=0\) se utiliza la construcción base.
Supongamos \(d\ge1\) y pongamos \(m=e(H)\). Se asigna un triángulo del
núcleo a cada arista exterior. Se exige que los triángulos sean
disjuntos en aristas y que sean disjuntos en vértices cuando sus aristas
exteriores sean incidentes. Antes de procesar \(xy\), los usos
anteriores en sus extremos excluyen a lo sumo \(6(d-1)\) vértices del
núcleo. Quedan \(k\ge p-6(d-1)\ge3p/4+6\) candidatos. Se han utilizado a
lo sumo \(3(m-1)\) aristas del núcleo. Como
\(m\le qd/2\le p^2/24+p/48\), \[
\begin{aligned}
\binom k2-3(m-1)-\left\lfloor k^2/4\right\rfloor
&\ge k^2/4-k/2-3m+3\\
&\ge p^2/64+29p/16+9>0.
\end{aligned} \tag{24}
\] La cota de Mantel proporciona un triángulo disponible. Para completar
el argumento, en un grafo sin triángulos sobre \(k\) vértices, la suma
de los grados de los extremos de cada arista es a lo sumo \(k\). Por
tanto, \(\sum_v d(v)^2\le ke\), mientras que Cauchy--Schwarz da
\(4e^2/k\le\sum_v d(v)^2\). Para \(e>0\) se obtiene \(e\le k^2/4\), y el
caso \(e=0\) es inmediato. Esto justifica la selección de (24) y
completa la inducción que construye los \(m\) bloques \(K_5\).

Sea \(t_a\) el número de estos bloques que contienen a un vértice \(a\)
del núcleo, y sea \(J\) el núcleo residual. Cada participación utiliza
dos aristas del núcleo y dos aristas cruzadas en \(a\). Por tanto, \[
d_J(a)=p-1-2t_a,\qquad |F_a|=2t_a,
\qquad d_J(a)+|F_a|=p-1. \tag{25}
\] Ahora se colorea \(E(J)\) de forma voraz con colores exteriores. Al
colorear \(ab\), las exclusiones de listas y las aristas incidentes
anteriores prohíben a lo sumo \[
|F_a|+|F_b|+(d_J(a)-1)+(d_J(b)-1)=2p-4
\] colores. Hay \(q\ge2p-1\) colores, por lo que siempre es posible
elegir uno. Se completa el núcleo con triángulos y se toman
individualmente las aristas cruzadas restantes. Los recuentos son \[
N_5=m,\qquad N_3=\binom p2-3m,\qquad N_2=pq-2\binom p2.
\] Su suma es \(M(n)-2m\), que alcanza la cota inferior. \(\square\)

La fórmula del valor por sí sola ya está cubierta por el Teorema A bajo
la condición menos restrictiva \(p\ge18d+1\). El contenido adicional
aquí es que toda arista exterior puede utilizar un \(K_5\) bajo
\(p\ge24d\), en lugar del umbral general de prescripción \(63d\), y que
la completación es voraz y no necesita el teorema HJ.

Aquí los colores cruzados que gasta un \(K_5\) se compensan exactamente
con la disminución del grado residual del núcleo. Esto permite una
completación voraz sin el teorema de coloreo con listas utilizado en el
caso general de perturbaciones.

\subsection{Apéndice B. Aristas en bloques
ineficientes}\label{apuxe9ndice-b.-aristas-en-bloques-ineficientes}

Bajo A, sea \(\Pi\) una partición con \(|\Pi|=\operatorname{cp}(G)+t\).
La ecuación (11) da \(t=\sum_{Q\in\Pi}\ell(a_Q,b_Q)\). A lo sumo \(t\)
bloques tienen holgura positiva, y la unión de sus aristas tiene tamaño
a lo sumo \[
21t. \tag{26}
\] Para demostrar esta última cota, pongamos \(k=a+b\). Si \(k\le7\) y
\(\ell(a,b)>0\), la integralidad da \(\binom k2\le21\le21\ell(a,b)\). Si
\(k\ge8\), se utiliza \[
40\ell(a,b)=(10b-4k-1)^2+4k^2-28k+39
\] para obtener \[
21\ell(a,b)-\binom k2
\ge\frac{(k-8)(64k-56)+371}{40}>0.
\] Al sumar se demuestra (26). La constante es óptima para esta
desigualdad por bloque: el tipo \((4,3)\) tiene holgura uno y 21
aristas. La cota localiza los bloques ineficientes; no implica una
reparación que deje fijos todos los demás bloques.

\subsection{Apéndice C. Muchas completaciones para cada
prescripción}\label{apuxe9ndice-c.-muchas-completaciones-para-cada-prescripciuxf3n}

Se fijan una prescripción y un levantamiento exterior del Teorema C con
carga a lo sumo \(7h\). Sea \(J\) el núcleo residual y sea \(L_e\) su
lista residual en una arista \(e\). Cada vértice del núcleo ha perdido a
lo sumo \(h\) vecinos originales del núcleo y a lo sumo \(14h\) aristas
del núcleo por el levantamiento. Por tanto, \[
\delta(J)\ge p-1-15h,\qquad
e(J)\ge\frac{p(p-1-15h)}2.
\] Su conjunto exterior prohibido tiene tamaño a lo sumo \(15h\), de
modo que \(|L_e|\ge q-30h\).

Pongamos \(d=\lfloor p/8\rfloor\). El teorema de Vizing {[}Viz64{]}
colorea \(J\) de forma propia con a lo sumo \(\Delta(J)+1\le p\)
colores. Se agrupan colores consecutivos en ventanas de a lo sumo \(d\)
colores. Hay a lo sumo \(\lfloor\Delta(J)/d\rfloor+1\le12\) ventanas
para \(p\ge16\). Una ventana de tamaño máximo contiene un subgrafo con
conjunto de aristas \(K\), grado máximo a lo sumo \(d\) y \[
|K|\ge e(J)/12\ge p^2/64.
\] Para la última desigualdad, \(p\ge63h\) y \(p\ge16\) implican
\(p-1-15h\ge3p/8\).

Se colorean las aristas de \(K\) en un orden fijo a partir de sus listas
residuales. Una arista tiene a lo sumo \(2(d-1)\) competidoras
incidentes en \(K\), y las hipótesis de tamaño dan \[
q-30h-2(d-1)\ge p.
\] Cada paso tiene al menos \(p\) elecciones, independientemente de las
elecciones legales anteriores. Esto cuenta al menos \(p^{|K|}\) prefijos
de coloreo propio con listas.

Para cualquiera de esos prefijos, se eliminan sus colores incidentes de
las listas de las aristas del núcleo no coloreadas. Cada lista pierde a
lo sumo \(2d\) colores adicionales y sigue teniendo tamaño al menos
\(q-30h-2d\ge p\). Se aplica Häggkvist--Janssen como en la Sección 2
para completar el núcleo residual conservando el prefijo. Dos prefijos
producen particiones distintas: para cada arista de \(K\), el único
triángulo que contiene esa arista del núcleo permite recuperar su color
exterior. El levantamiento exterior permanece fijo en todo el proceso,
por lo que ninguno de sus bloques utiliza una arista de \(K\). Esto
demuestra la cota inferior para una sola prescripción.

Las familias correspondientes a prescripciones distintas son disjuntas,
puesto que el orden del único bloque que pasa por una arista exterior
permite recuperar su etiqueta. Multiplicar por \(2^m\) da la cota
inferior total.

Para la cota superior, todo bloque óptimo tiene orden a lo sumo cinco,
por la identidad de holgura. Hay a lo sumo \((n+1)^5\) conjuntos de
vértices de tamaño a lo sumo cinco: se ordenan sus elementos de una
manera fija y se rellenan las posiciones restantes con un símbolo
centinela para obtener una codificación inyectiva mediante cinco
entradas de un alfabeto de tamaño \(n+1\). Una partición está
determinada por la aplicación que asigna a cada arista su único bloque.
Por tanto, hay a lo sumo \(((n+1)^5)^{e(G)}\) particiones óptimas.

Por último, \(n=\Theta(p)\), \(e(G)=O(p^2)\) y el exponente inferior es
\(\Theta(p^2)\). Al tomar logaritmos se obtiene el orden de crecimiento
enunciado. El complemento formal registra las cotas finitas explícitas,
sin exigir como premisa un enunciado asintótico separado. \(\square\)

\subsection{Agradecimientos}\label{agradecimientos}

El autor agradece profundamente a su esposa María Paz y a sus hijos
Lucas, Juan Cristóbal, Francisca, Raimundo y Benjamín su amor, paciencia
y apoyo.

\subsection{Uso de inteligencia artificial y herramientas
computacionales}\label{uso-de-inteligencia-artificial-y-herramientas-computacionales}

Se utilizó ChatGPT/Codex, de OpenAI, para explorar argumentos,
comprobarlos y preparar el manuscrito. Aristotle, de Harmonic, participó
en la exploración de caminos, la búsqueda de contraejemplos, la
elaboración de pruebas y su formalización y revisión en Lean; no sólo
tradujo pruebas terminadas. El Paper IV acredita a Claude, de Anthropic,
por su contribución al trabajo heredado.

Certo {[}Certo{]} se utilizó en las nuevas comprobaciones finitas. Las
comprobaciones anteriores realizadas con Certo y Jacobian {[}Jacobian{]}
se conservan en los materiales de auditoría adjuntos, donde se
identifican sus versiones, especificaciones y resultados registrados.
Los certificados finitos se comprueban por separado de la construcción
que los produce. Su uso se declara por separado de las dependencias
lógicas de la prueba formal. Los cálculos se revisaron según su alcance;
una búsqueda finita no reemplaza una prueba universal.

El expediente de fuentes registra dos revisiones de la construcción por
IA, una de ellas con participación en las derivaciones, y una auditoría
algebraica. Estas no son revisiones humanas independientes.

El autor conserva la responsabilidad por los argumentos, las citas, el
código y la presentación. Ningún sistema de IA figura como autor.

\subsection{Referencias}\label{referencias}

{[}EFO88{]} P. Erdős, R. J. Faudree and E. T. Ordman, \emph{Clique
partitions and clique coverings}, Discrete Mathematics \textbf{72}
(1988), 93--101.
\href{https://www.renyi.hu/~p_erdos/1988-04.pdf}{Original article}.
\href{https://doi.org/10.1016/0012-365X(88)90197-5}{DOI}.

{[}HJ97{]} R. Häggkvist and J. C. M. Janssen, \emph{New Bounds on the
List-Chromatic Index of the Complete Graph and Other Simple Graphs},
Combinatorics, Probability and Computing \textbf{6} (1997), 295--313.
\href{https://doi.org/10.1017/S0963548397002927}{DOI}.

{[}Mon96{]} S. D. Monson, \emph{The effects of vertex deletion and edge
deletion on the clique partition number}, Ars Combinatoria \textbf{42}
(1996), 89--96.
\href{https://combinatorialpress.com/article/ars/Volume\%20042/volume-42-paper-6.pdf}{Original
article}.

{[}Oke26{]} O. Okechukwu, \emph{Clique partitions and bounded simplicial
defect}, arXiv:2609.20871v1 (2026), especially Lemmas 4.1--4.3.
\href{https://arxiv.org/html/2609.20871v1}{Preprint}.

{[}Nin26{]} B. Ning, \emph{On the difference between clique partition
and clique covering numbers of graphs}, arXiv:2608.11536v1 (2026),
Proposition 3.1. \href{https://arxiv.org/html/2608.11536v1}{Preprint}.

{[}Viz64{]} V. G. Vizing, \emph{On an estimate of the chromatic class of
a p-graph}, Diskret. Analiz. \textbf{3} (1964), 25--30 (in Russian).

{[}IV{]} J. P. Traverso Gianini, \emph{Clique partitions with rooted
simplicial defect: quantitative stability and sharp eventual bounds},
Paper IV, preprint v1.23 (2026).

{[}V{]} J. P. Traverso Gianini, \emph{Colorability and Realizability
Defects in Resource-Dominated Packing Relaxations}, Paper V, Preprint
1.0 (2026).
\href{https://github.com/jtraverso/erdos-81-chordal-clique-partitions/tree/main/preprints/PAPER_V}{Fuentes
y manuscrito públicos}.

{[}PD81{]} N. J. Pullman and A. Donald, \emph{Clique coverings of
graphs---II: complements of cliques}, Utilitas Mathematica \textbf{19}
(1981), 207--213. Citado mediante {[}EFO88, p.~94{]} y {[}WZ93,
Corolario 3.3{]}; el artículo original no se consultó para esta
revisión.

{[}EOZ93{]} P. Erdős, E. T. Ordman and Y. Zalcstein, \emph{Clique
partitions of chordal graphs}, Combinatorics, Probability and Computing
\textbf{2} (1993), 409--415.
\href{https://doi.org/10.1017/S0963548300000808}{DOI}.

{[}CEO94{]} G.-T. Chen, P. Erdős and E. T. Ordman, \emph{Clique
partitions of split graphs}, in Y. Alavi, D. R. Lick and J. Liu (eds.),
\emph{Combinatorics, Graph Theory, Algorithms and Applications}
(Beijing, 1993), World Scientific, 1994, 21--30.
\href{http://ordman.net/MathResearch/CEOClique_Parts.pdf}{Author's
copy}.

{[}WW91{]} W. D. Wallis and J. Wu, \emph{On clique partitions of split
graphs}, Discrete Mathematics \textbf{92} (1991), 427--429.
\href{https://doi.org/10.1016/0012-365X(91)90297-F}{DOI}.

{[}WZ93{]} W. D. Wallis and G.-H. Zhang, \emph{On the partition and
coloring of a graph by cliques}, Discrete Mathematics \textbf{120}
(1993), 191--203.
\href{https://doi.org/10.1016/0012-365X(93)90576-F}{DOI}.

{[}Ma88{]} S.-H. Ma, \emph{Clique partitions of threshold graphs}, Acta
Mathematicae Applicatae Sinica \textbf{11} (1988), 493--502.
\href{https://applmath.cjoe.ac.cn/jweb_yysxxb/EN/10.12387/C1988054}{Publisher}.
En chino, con resumen en inglés.

{[}MR90{]} S. McGuinness y R. Rees, \emph{On the number of distinct
minimal clique partitions and clique covers of a line graph}, Discrete
Mathematics \textbf{83} (1990), 49--62.
\href{https://faculty.tru.ca/smcguinness/cliquecoverslinegraph.pdf}{Copia
del autor}. \href{https://doi.org/10.1016/0012-365X(90)90220-C}{DOI}.

{[}KSS23{]} M. Kwan, A. Sah y M. Sawhney, \emph{Enumerating Matroids and
Linear Spaces}, Comptes Rendus. Mathématique \textbf{361} (2023),
565--575.
\href{https://numdam.org/articles/10.5802/crmath.423/}{Artículo
publicado}.

{[}IKD12{]} T. Ito, M. Kamiński y E. D. Demaine, \emph{Reconfiguration
of List Edge-Colorings in a Graph}, Discrete Applied Mathematics
\textbf{160} (2012), 2199--2207.
\href{https://doi.org/10.1016/j.dam.2012.05.014}{DOI}. La
\href{https://erikdemaine.org/papers/NPReconfiguration_WADS2009/paper.pdf}{versión
WADS 2009} está disponible en el sitio de los autores.

{[}HJ-formal{]} J. P. Traverso Gianini, \emph{The list chromatic index
of complete graphs (Häggkvist--Janssen)}, Lean formalization, Palomar
\href{https://palomar-registry.org/entry.html?id=PALOMAR-2026-10-15-000003&version=1}{PALOMAR-2026-10-15-000003,
version 1}.

{[}Viz-formal{]} J. P. Traverso Gianini, \emph{Vizing's theorem and
equitable edge colourings}, Lean formalization, Palomar
\href{https://palomar-registry.org/entry.html?id=PALOMAR-2026-10-15-000004&version=1}{PALOMAR-2026-10-15-000004,
version 1}.

{[}VI-companion{]} J. P. Traverso Gianini, \emph{Formalization and
reproducibility companion to Optimal Clique Partitions under Local
Perturbations}, accompanying source distribution, Preprint 1.0 (2026).
\href{https://github.com/jtraverso/erdos-81-chordal-clique-partitions/tree/main/preprints/PAPER_VI/05_formalization}{Fuentes
y complemento públicos}. Exact source and archive hashes are recorded in
the companion.

{[}Certo{]} J. P. Traverso Gianini, \emph{Certo: computational tools
with verifiable certificates}.
\href{https://github.com/jtraverso/certo-math}{Software}. The audited
local version and certificate specifications are fixed in the companion.

{[}Jacobian{]} Morluto and contributors, \emph{Jacobian: composable
mathematical tools}.
\href{https://github.com/morluto/jacobian}{Software}. Installed
operation declarations and versions used here are recorded with the
computational checks.

\end{document}
